\documentclass[11pt]{article}
\usepackage[margin=1in]{geometry}
\usepackage{amsmath}
\usepackage{booktabs}
\usepackage{longtable}
\title{Replacement memo: corrected AR(1) Section 5.3}
\author{SVP simulation audit}
\date{15 September 2026}
\begin{document}
\maketitle

\section*{Decision}
Section 5.3 is replaced by the calibrated four-method AR(1) power study
described below. The original \texttt{SVP\_NEW.pdf} remains untouched because
no editable source was available. The pre-refresh AR(1) outputs are preserved
in \texttt{simulations/other\_simus/power\_ar1\_legacy\_20260915/}.

\section*{Study design}
We simulate independent realizations of the piecewise-constant mean model
\[
Y_t = \mu_t + E_t, \qquad E_t = \rho E_{t-1} + \eta_t,
\qquad \rho=0.8,
\]
with $n=600$ observations and unit stationary marginal noise variance. The
innovations are Gaussian with variance $1-\rho^2$, and the first residual is
drawn from the stationary distribution. The four signal patterns are
\texttt{none}, \texttt{up}, \texttt{updown}, and \texttt{rand1}. Jump sizes
range from 1 to 5 in increments of 0.2, with 100 replications per
configuration.

\section*{Methods and calibration}
The competitors are PELT AR1 approximate, PELT inflated, DeCAFS AR1, and SVP
AR1Focus. The approximate PELT method applies PELT to the transformed series
\[
Z_t=Y_{t+1}-\rho Y_t,
\]
and maps each transformed endpoint $j$ to endpoint $j+1$ in the original
series. It is approximate because ordinary transformed-data least squares does
not retain the complete AR(1) likelihood boundary correction.

DeCAFS uses the pure AR(1) state-space model with
\texttt{sdEta = 0}, \texttt{sdNu = sqrt(1-rho\string^2)}, and
\texttt{phi = rho}. Its multiplier $c_D$ in $\beta=c_D\log(n)$ is selected as
the smallest candidate attaining corrected null F1 at least 0.99. The
approximate PELT multiplier and the AR1Focus threshold are selected to match
the selected DeCAFS null F1. AR1Focus uses \texttt{subtests = "right"} only.
The inflated PELT baseline retains the fixed penalty
$3\log(n)(1+\rho)/(1-\rho)$. The approximate-Pelt multiplier is calibrated
alongside DeCAFS and AR1Focus.

Calibration and independent validation each use 1,000 stationary AR(1) null
series, deterministic seeds, and Wilson confidence intervals. The complete
candidate table and selected values are stored with the study outputs.

\begin{table}

\caption{Validation null F1 for the selected AR(1) methods.}
\centering
\begin{tabular}[t]{lrrrr}
\toprule
method & candidate & null\_f1 & null\_f1\_lower & null\_f1\_upper\\
\midrule
DeCAFS AR1 & 7.75 & 0.985 & 0.9754 & 0.9909\\
PELT AR1 approximate & 0.80 & 0.989 & 0.9804 & 0.9938\\
SVP AR1Focus / right & 1.85 & 0.993 & 0.9856 & 0.9966\\
PELT inflated & 3.00 & 0.999 & 0.9944 & 0.9998\\
\bottomrule
\end{tabular}
\end{table}


\section*{Metrics}
For every method, the mandatory terminal endpoint $n$ is excluded before
one-to-one matching within a tolerance of five observations. We report
precision, recall, F1, the probability of recovering the correct number of
changepoints, signal MSE, the number of fitted segments, and localization
error. The same metric implementation is used by the Gaussian and AR(1)
studies.

\section*{Before/after comparison}
The detailed paired comparison is stored in
\texttt{simulations/z\_AR1/paired\_comparison.csv}; it compares the refreshed
results with the archived pre-refresh study by pattern and jump size.

{\scriptsize
\setlength{\tabcolsep}{2pt}
\begin{table}

\caption{Mean refreshed-minus-legacy metric change by pattern and method.}
\centering
\begin{tabular}[t]{llrrrrrr}
\toprule
pattern & method & F1 & Precision & Recall & Correct & MSE & Segments\\
\midrule
none & DeCAFS AR1 & 0.0381 & 0.0381 & 0.0381 & 0.0381 & -0.0016 & -0.0381\\
none & PELT AR1 approximate & -0.0090 & -0.0090 & -0.0090 & -0.0090 & 0.0016 & 0.0105\\
none & PELT inflated & 0.0000 & 0.0000 & 0.0000 & 0.0000 & 0.0000 & 0.0000\\
none & SVP AR1Focus & -0.0033 & -0.0033 & -0.0033 & -0.0033 & 0.0004 & 0.0033\\
rand1 & DeCAFS AR1 & -0.0797 & -0.0704 & -0.0731 & -0.0019 & 0.1612 & -0.5376\\
\addlinespace
rand1 & PELT AR1 approximate & 0.1037 & 0.0427 & 0.1635 & -0.0200 & -0.3443 & 2.7629\\
rand1 & PELT inflated & 0.0000 & 0.0000 & 0.0000 & 0.0000 & 0.0000 & 0.0000\\
rand1 & SVP AR1Focus & 0.1177 & 0.0932 & 0.1168 & 0.0143 & -0.2183 & 0.9067\\
up & DeCAFS AR1 & -0.0477 & 0.0039 & -0.0689 & -0.1290 & 0.2085 & -0.5633\\
up & PELT AR1 approximate & -0.0039 & -0.1705 & 0.1678 & -0.0252 & -0.6099 & 4.0681\\
\addlinespace
up & PELT inflated & 0.0000 & 0.0000 & 0.0000 & 0.0000 & 0.0000 & 0.0000\\
up & SVP AR1Focus & 0.1367 & 0.0652 & 0.1605 & 0.2305 & -0.5985 & 1.0838\\
updown & DeCAFS AR1 & -0.1578 & -0.1365 & -0.1599 & -0.1438 & 0.3199 & -1.1590\\
updown & PELT AR1 approximate & 0.2168 & 0.1307 & 0.3476 & -0.0210 & -0.7302 & 5.4981\\
updown & PELT inflated & 0.0000 & 0.0000 & 0.0000 & 0.0000 & 0.0000 & 0.0000\\
\addlinespace
updown & SVP AR1Focus & 0.2416 & 0.2467 & 0.2394 & 0.2395 & -0.4443 & 1.7462\\
\bottomrule
\end{tabular}
\end{table}

}

\end{document}
